\documentclass[10pt,landscape,a4paper]{article}
\usepackage[utf8]{inputenc}
\usepackage[UKenglish]{babel}
\usepackage{tikz}
\usetikzlibrary{shapes,positioning,arrows,fit,calc,graphs,graphs.standard}
\usepackage[nosf]{kpfonts}
\usepackage[t1]{sourcesanspro}
%\usepackage[lf]{MyriadPro}
%\usepackage[lf,minionint]{MinionPro}
\usepackage{multicol}
\usepackage{wrapfig}
\usepackage[top=0mm,bottom=4mm,left=2mm,right=1mm]{geometry}
\usepackage[framemethod=tikz]{mdframed}
\usepackage{microtype}
\usepackage{mathtools} % loads amsmath, for math environments etc
\usepackage{setspace}
\usepackage{tikz}
\usepackage{tikz}
\usepackage{enumitem}

\let\bar\overline

\definecolor{myblue}{cmyk}{1,.72,0,.38}

\def\firstcircle{(0,0) circle (1.5cm)}
\def\secondcircle{(0:2cm) circle (1.5cm)}

\colorlet{circle edge}{myblue}
\colorlet{circle area}{myblue!5}

\tikzset{filled/.style={fill=circle area, draw=circle edge, thick},
    outline/.style={draw=circle edge, thick}}

\pgfdeclarelayer{background}
\pgfsetlayers{background,main}

\everymath\expandafter{\the\everymath \color{myblue}}
\everydisplay\expandafter{\the\everydisplay \color{myblue}}

\renewcommand{\baselinestretch}{.8}
\pagestyle{empty}

\global\mdfdefinestyle{header}{%
linecolor=gray,linewidth=1pt,%
leftmargin=6mm,rightmargin=6mm,skipbelow=1mm,skipabove=1mm,
}

\newcommand{\header}{
\begin{mdframed}[style=header]
\scriptsize

\textbf{Mosteller Cheat sheet by~Alain~Chenier,~page~\thepage~of~2,  \today{} }
\end{mdframed}
}

\makeatletter
\renewcommand{\section}{\@startsection{section}{1}{0mm}%
                                {.1ex}%
                                {.1ex}%x
                                {\color{blue}\sffamily\small\bfseries}}
\renewcommand{\subsection}{\@startsection{subsection}{1}{0mm}%
                                {.1ex}%
                                {.1ex}%x
                                {\sffamily\bfseries}}



\def\multi@column@out{%
   \ifnum\outputpenalty <-\@M
   \speci@ls \else
   \ifvoid\colbreak@box\else
     \mult@info\@ne{Re-adding forced
               break(s) for splitting}%
     \setbox\@cclv\vbox{%
        \unvbox\colbreak@box
        \penalty-\@Mv\unvbox\@cclv}%
   \fi
   \splittopskip\topskip
   \splitmaxdepth\maxdepth
   \dimen@\@colroom
   \divide\skip\footins\col@number
   \ifvoid\footins \else
      \leave@mult@footins
   \fi
   \let\ifshr@kingsaved\ifshr@king
   \ifvbox \@kludgeins
     \advance \dimen@ -\ht\@kludgeins
     \ifdim \wd\@kludgeins>\z@
        \shr@nkingtrue
     \fi
   \fi
   \process@cols\mult@gfirstbox{%
%%%%% START CHANGE
\ifnum\count@=\numexpr\mult@rightbox+2\relax
          \setbox\count@\vsplit\@cclv to \dimexpr \dimen@-1cm\relax
\setbox\count@\vbox to \dimen@{\vbox to 1cm{\header}\unvbox\count@\vss}%
\else
      \setbox\count@\vsplit\@cclv to \dimen@
\fi
%%%%% END CHANGE
            \set@keptmarks
            \setbox\count@
                 \vbox to\dimen@
                  {\unvbox\count@
                   \remove@discardable@items
                   \ifshr@nking\vfill\fi}%
           }%
   \setbox\mult@rightbox
       \vsplit\@cclv to\dimen@
   \set@keptmarks
   \setbox\mult@rightbox\vbox to\dimen@
          {\unvbox\mult@rightbox
           \remove@discardable@items
           \ifshr@nking\vfill\fi}%
   \let\ifshr@king\ifshr@kingsaved
   \ifvoid\@cclv \else
       \unvbox\@cclv
       \ifnum\outputpenalty=\@M
       \else
          \penalty\outputpenalty
       \fi
       \ifvoid\footins\else
         \PackageWarning{multicol}%
          {I moved some lines to
           the next page.\MessageBreak
           Footnotes on page
           \thepage\space might be wrong}%
       \fi
       \ifnum \c@tracingmulticols>\thr@@
                    \hrule\allowbreak \fi
   \fi
   \ifx\@empty\kept@firstmark
      \let\firstmark\kept@topmark
      \let\botmark\kept@topmark
   \else
      \let\firstmark\kept@firstmark
      \let\botmark\kept@botmark
   \fi
   \let\topmark\kept@topmark
   \mult@info\tw@
        {Use kept top mark:\MessageBreak
          \meaning\kept@topmark
         \MessageBreak
         Use kept first mark:\MessageBreak
          \meaning\kept@firstmark
        \MessageBreak
         Use kept bot mark:\MessageBreak
          \meaning\kept@botmark
        \MessageBreak
         Produce first mark:\MessageBreak
          \meaning\firstmark
        \MessageBreak
        Produce bot mark:\MessageBreak
          \meaning\botmark
         \@gobbletwo}%
   \setbox\@cclv\vbox{\unvbox\partial@page
                      \page@sofar}%
   \@makecol\@outputpage
     \global\let\kept@topmark\botmark
     \global\let\kept@firstmark\@empty
     \global\let\kept@botmark\@empty
     \mult@info\tw@
        {(Re)Init top mark:\MessageBreak
         \meaning\kept@topmark
         \@gobbletwo}%
   \global\@colroom\@colht
   \global \@mparbottom \z@
   \process@deferreds
   \@whilesw\if@fcolmade\fi{\@outputpage
      \global\@colroom\@colht
      \process@deferreds}%
   \mult@info\@ne
     {Colroom:\MessageBreak
      \the\@colht\space
              after float space removed
              = \the\@colroom \@gobble}%
    \set@mult@vsize \global
  \fi}

\makeatother
\setlength{\parindent}{0pt}
\setlength{\columnseprule}{0.01pt}

\begin{document}
\scriptsize
\begin{multicols*}{3}
	
\section{Sock Drawer}
\subsection*{Question}
drawer contains r red sock and b blue socks. 2 socks are drawn. \textbf{force probability to be 1/2}.\\
a) how small num of red and black socks\\
b) how small num of socks if num black socks is even. 

\subsection*{Answer}
\begin{spacing}{0.1}

$$\text{prob first is red}\frac{r}{r+b} \text{ then prob second is red given first is red}\frac{r-1}{r+b-1}$$ 
$$\text{so need to find r b such that}  \left( \frac{r}{r+b}\right) \left(  \frac{r-1}{r+b-1}\right)  = \frac{1}{2}  $$
$$\frac{r}{r+b} > \frac{r-1}{r+b-1} \text{so with 1/2 condition above}$$
$$\frac{r}{r+b} > \frac{1}{\sqrt{2}} > \frac{r-1}{r+b-1} \text{ so } (\sqrt{2}+1)b+1 > r >  (\sqrt{2}+1)b$$
\end{spacing}

\section{Win 2 tennis matches in a row}
\subsection*{Question}
win at least 2 tennis matches in a row in a 3 set with Champion > Father\\
\textbf{Q:} Choose Father-Champion-Father or Champion - Father - champion ? 
\subsection*{Answer}
Compare both sequences FCF CFC, add up the probs of winning:\\

\begin{tabular}{|c|c|c||c|}
	\hline 
	\textbf{F} &\textbf{ F} & \textbf{F} & \textbf{Prob} \\ 
	\hline 
	W & W & W & fcf \\ 
	W & W & L & fc(1-f) \\ 
	L & W & W & (1-f)cf \\ 
	\hline 
	\textbf{Total} &  &  & \textbf{fc (2-f)} \\ 
	\hline 
\end{tabular} 
\quad
\begin{tabular}{|c|c|c||c|}
	\hline 
    \textbf{F} &\textbf{ F} & \textbf{F} & \textbf{Prob} \\ 
	\hline 
	W & W & W & cfc \\ 
	W & W & L & cf(1-c) \\ 
	L & W & W & (1-c)fc \\ 
	\hline 
	\textbf{Total} &  &  & \textbf{fc(2-c)} \\ 
	\hline 
\end{tabular} 
\\

\section{Jury: 2 with  $p$ + random or 1 with $p$ ?}
\subsection*{Answer}
prepare table of all correct decisions:\\
\begin{tabular}{c|c|c|c||c|}
	\hline 
	Serious agree & p & p & not matter as overruled & $p^2$ \\ 
	Serious disagree 1 & p & 1-p & chooses right with 1/2 & $p(1-p) \frac{1}{2}$ \\ 
	Serious disagree 2 & 1-p & p & chooses right with 1/2 & $(1-p)p \frac{1}{2}$\\ 
	\hline 
	\textbf{Total} &  &  &  & $p^2 + p(1-p) = p$\\ 
	\hline 
\end{tabular}

\section{How many throws of dice to get 6 ?}
\subsection*{Answer 1 - force}
\begin{spacing}{0.1}
let $p$ = prob of 6 = 1/6 then prob of success $p$ (len=1). $qp$ (len2), $qqp$ (len 3) ...\\
$$m=1(p) + 2(qp) + 3 (qqp)\text{...}=p(1+ 2q+ 3qq + \text{...})  = p \left( \dfrac{1}{1-q}^{2} \right)= p \dfrac{1}{p^2}=\frac{1}{p}$$

\subsection*{Answer 2 - First step analysis = Markov FSA}
If first is failure, then m->m+1, if first is good then m=1 so\\
$$m=q(m+1)+p(1) \Leftrightarrow m(1-q)=q+p=1 \Leftrightarrow m=\dfrac{1}{p}$$

\subsection*{Answer 3 - Martingale}
...put martingale here...

\end{spacing}

\begin{spacing}{0.1}

\section{throw 3/4cm diameter coin on 1cm x 1cm grided table}
\subsection*{Question}
Get 5c (in all) if whole inside a square , lose 1c otherwise. CHANCE TO WIN and How MUCH on AVG?
\subsection*{Answer}
Centre of coin falls in one square, so consider one square, and divide allowable area for centre by area of square\\

\begin{tabular}{c|c}
\begin{tikzpicture} [scale=0.30]
\draw (0,0) rectangle (1,1);
\draw (0.4,0.4) circle (0.375cm);
\draw [fill=orange] (0, 0) + (0.5, 0.5) rectangle (0.75,0.75);
\end{tikzpicture}
&
$p = \frac{1}{4} \text{x}\frac{1}{4} \text{and } m= (+5) \frac{1}{16} + (-1) \frac{15}{16}$ 
\end{tabular} 

\end{spacing}

\begin{spacing}{0.01}
\section{Dice Game}
\subsection*{Question}
Bet on 1 or 2 or ... or 6 and throw 3 dice. Receive 1x-2x-3x stake + original stake. WHAT IS the EXPECTED LOSS ? 
\subsection*{Answer}
when you roll either all 3 diff, or 2 same, or 3 same. Bet 6 on all 6. If all 3 different then house wins 3, you lose 3 (tot 0), if 2 same h wins 1(for die=3)+1 (for die=4)+1(for 5)+1(for 6), you win 2(for 1)+1(for 2) (tot you=3-4=\textbf{-1}). if 3 same you win 3 house wins 5 (tot you = 3-5=\textbf{-2}). Then count how many  ways to get 1 alike, 2 alike, 3 alike and multiply by loss.\\
Total ways = $6x6x6=216$, all 3 different = $6x5x4=120$, all 3 same = $6x1x1=6$, therefore 2 same $216-120-6 =216-126=90$ [Note: 2 same (alt way) : die 1=2 and 3 dif). or (die 1=3 and 2 dif) or (die 2=3 and 1 dif)so $6x5 + 6x5 + 6x5 = 30 x 3 =90 $] So loss for betting on all 6 = all diff + 2 same + 3 same = $ \frac{0}{6} \frac{120}{216} + \frac{(-1)}{6} \frac{90}{216} + \frac{(-2)}{6} \frac{6}{216} = \frac{-102}{216} = \frac{-17}{16} $ 
\end{spacing}

\begin{spacing}{0.1}

\section{Follow-up Bet on outcome of initial roulette bet}
\subsection*{Question}
Bet 1 number out of 38 then bet 20 ahead after 36 games. HOW MUCH IS THE WIN or LOSS ?
\subsection*{Answer}
$(\text{expected \textbf{gain} one turn}) = (+35) \frac{1}{38} + (-1)\frac{37}{38} = \frac{-2}{38}$ , and in 36 turns $ (\text{expected \textbf{gain} 36 turns}) = m_1 = (36) \frac{-2}{38} = -\frac{72}{38} $
\\
Now for the 20 bet: 'ahead in 36 games' = 'win once or more in 36 games' since \textbf{gain} (for win 0 times in 36) = $36(-1)=-36 < 0$ , and \textbf{gain} (win 1 in 36) = $1(+35) + 35(-1)=0 >= 0$
\\
so relevant prob for the 20 bet = P(win one or more in 36)= 1- P(lose all 36) = $1-\left( \frac{37}{38}\right) ^{36} = 0.617$ so for the 20 bet expected \textbf{gain} = $m_2=(+20) (+0.617) + (-20) (1-0.617) $ 
\\
and so total \textbf{gain} $m = m_1 + m_2 = ...$

\section{All cards same suit in 52 card pack}
\subsection*{Question}
Probability of all 13 cards of same suit?
\subsection*{Answer 1 - straight}
$P=\frac{\text{good hands}}{\text{all hands}}$. for good: pick any card ,  after must be of same suit = $ 52 * 12 * 11 * ... * 1 = 52 * 12!$, for all hands: any card then, any card left, ... (13 times)= $52 * 51 * ... * 40 = \frac{52!}{39!}$ so Prob = $\frac{52 * 12!}{\frac{52!}{39!}} =\frac{39!12!}{51!} $ \\
\subsection*{Answer 2 - binomials}
$P=\frac{\text{good hands}}{\text{all hands}}$. for good: pick 1 of 4 suits, then pick 13 cards in the 13 card suits , for all, pick 13 cards in 52, so $P=\frac{\binom{4}{1} \binom{13}{13} }{\binom{52}{13}} = \frac{(4)(1)}{\frac{52!}{39!13!}} = \frac{4*13!*39!}{52!} = \frac{4*13*12!*39!}{52 * 51!}=\frac{12!*39!}{51!}$= same as above!!
\section{Craps game}
\subsection*{Game description and question}
Only total of sum of 2 dice counts. Lose if sum=2,3,12. Win if sum=7,11. record sum otherwise, then keep throwing until either 'same sum again' (win) or  7 (lose). \textbf{Question}: odds of winning ? 
\subsection*{Answer 1 - Conditional probability}
get probs of totals by counting in the matrix:
\begin{tabular}{|c||c|c|c|c|c|c|}
	\hline 
	& \textbf{\textit{1}} & \textbf{\textit{2}} & \textbf{\textit{3}} & \textbf{\textit{4}} & \textbf{\textit{5}} & \textbf{\textit{6}} \\ 
	\hline 
	\hline 
	\textbf{\textit{1}} & 2 & 3 & 4 & 5 & 6 & 7 \\ 
	\hline 
	\textbf{\textit{2}} & 3 & 4 & 5 & 6 & 7 & 8 \\ 
	\hline 
	\textbf{\textit{3}} & 4 & 5 & 6 & 7 & 8 & 9 \\ 
	\hline 
	\textbf{\textit{4}} & 5 & 6 & 7 & 8 & 9 & 10 \\ 
	\hline 
	\textbf{\textit{5}} & 6 & 7 & 8 & 9 & 10 & 11 \\ 
	\hline 
	\textbf{\textit{6}} & 7 & 8 & 9 & 10 & 11 & 12 \\ 
	\hline 
\end{tabular} so

\begin{tabular}{|c||c|c|c|c|c|c|c|c|c|c|c|}
	\hline 
	total & 2 & 3 & 4 & 5 & 6 & 7 & 8 & 9 & 10 & 11 & 12\\ 
	\hline 
	p(total) & $\frac{1}{36}$ & $\frac{2}{36}$ & $\frac{3}{36}$ & $\frac{4}{36}$ & $\frac{5}{36}$ & $\frac{6}{36}$ & $\frac{5}{36}$ & $\frac{4}{36}$ & $\frac{3}{36}$ & $\frac{2}{36}$ & $\frac{1}{36}$\\ 
	\hline 
\end{tabular} 

	so $P(\text{win 1st throw})=P(7)+P(11)=\frac{6}{36}+\frac{2}{36} \text{ and } P(\text{loss 1st throw})=P(2)+P(3)+P(12)=\frac{1}{36}+\frac{2}{36}+\frac{1}{36}$ then after first throw only: throwing either 'initial sum again' (which was either=4,5,6,8,9 or 10) (win) or 7 (lose) matters, so we have a conditional prob on first throw.  Say the 'initial sum was 4', then win if 4, or lose if 7, so P(win if init sum=4)=P(win|4)=$\frac{\text{num ways of 4}}{\text{num ways of 4 or 7}}=\frac{3}{3+6}$ so P(win if 4)=P(4)*P(win|4), so in total P(win)=P(7)+P(11)+P(4)*P(win|4)+P(5)*P(win|5)+P(6)*P(win|6)+P(7)*P(win|7) \\ 
	+P(8)*P(win|8)+P(9)*P(win|9)+P(10)*P(win|10) and expected gain = (+1)*P(win)+(-1)*P(loss)
\subsection*{note: Answer 1' - infinite summation method untill you get the init sum or 7}
if into 2nd or more throw , you have the possibilities P(init sum again), P(ignore ie not count), or P(7), and win if 'init again', or, 'not count'+'init again, or, 'not count' 'not count' then 'init again' etc. So say init sum was 4 (same as above), then P(win|init sum=4) = P+RP+RRP+... with P=prob(init sum of 4 again) and R=prob(not count), which is P(1+R+RR+..)=P/(1-R) , so in summary to win with a 4 then P(4)*P/(1-R) with P=3/36 and R=1-P(4)-P(7)=1-3/36-6/36 -- note P/(1-R)  = P/(P+Q) <1 if Q=prob(7)


\section{Which game do you prefer}
\subsection*{Question: prefer known contents or unknown contents?}
game1:urn1:10 black balls,urn2: 10 white balls, pick a colour
game2:urn1:unknown,urn2: urn2, pick a colour
\subsection*{Answer: same}
game1: P(win)=P(urn1)P(B|urn1)+P(urn2)P(B|urn1)=1/2 1/2 +  1/2 1/2 = 1/2 for B=black, same for W\\
game2: P(win)=P(urn1)P(B|urn1)+P(urn2)P(B|urn1)=1/2 x +  1/2 (1-x) = 1/2 for B=black, same for W because $x = \blacksquare/total$ so in urn2: $(total- \blacksquare)/total=1-x$


\section{Prisoner's dilemma}
\subsection*{Question - find the flaw in this argument}
3 prisoners A,B,C and 2 will be released. \textbf{W}arden knows. if A asks W, W will name 1 other than A. A thinks his chance before asking =2/3, and if \textbf{W}arden say 'B' then chance now = 1/2 < 2/3, hence should not ask. Where the flaw ? 
\subsection*{Answer: list ALL events, must condition on the Warden}
\begin{tabular}{|c|c|c|c|c|}
	\hline 
	& & &  &  \\ 
	\hline 
	Num & Release Outcome & Warden & Prob outcome AND W & example W says B \\ 
	\hline 
	1 & \textbf{A,B} & C must & (1/3)(1)=1/3 &  \\ 
	\hline 
	2 & \textbf{A,C} & B must & (1/3)(1)=1/3 & ok\\ 
	\hline 
	3 & B,C & B & $(1/3)(1/2)=1/6$ & ok \\ 
	\hline 
	4 & B,C & C & $(1/3)(1/2)=1/6$ &  \\ 
	\hline 
\end{tabular} 
so P(A released | W says B) = P(A released and W says B) /P(W says B)= row2/(row2+row3) = (1/3)/(1/6+1/3) = (1/3)/(3/6) = 6/9 = 2/3 = same as without asking
similarly P(A released | W says C) = P(A released and W says C) /P(W says C)= row1/(row1+row3) = same since same probs = same as without asking

\section{Coupon Collector}
\subsection*{Question: How many boxes required to collect all 5 coupons}
5 possible coupons 1..5, 1 per box, how many boxes on average to get all 5?
\subsection*{Answer: See FSA problem 4 - num throws to get a 6 = 1/(1/6)=6 }
1st box always good. second box:new number with prob (4/5) $\rightarrow$  from FSA prob 4 must wait 1/(4/5)=5/4, third box: new number with prob (3/5) $\rightarrow$ must wait 1/(3/5)=5/3 etc so answer $5/4+5/3+5/2+5/1=5(1/4+1/3+1/2+1)$. Note $1+1/2+..+1/n=\ln n + \frac{1}{2n} + euler$ so $n \ln n+ 1/2 +n \gamma \equiv n \ln n +n \gamma \text{for big n}$

\section{Theater Row - pairs of adjacent but different}
\subsection*{Question: How many pairs of adjacent but different}
8 B and 7 M in a row, how many pairs of \textbf{BM} or \textbf{MB} on average? For example \textbf{BMB}MMMMM\textbf{MB}BBB is 3 good pairs. \textbf{Then same question with 6B 5M 4W}
\subsection*{Answer1: prob 2 balls are adjacent x num of good pairs}
prob 2 balls in the 15 are adjacent = 14 slots possible / (15 choose 2) , then num pairs=8x7=56 so $ \frac{14}{\binom{15}{2}}(56) = \frac{14}{ \frac{15 x 14}{2}} (56) = (2) \frac{56}{15}$ . For next one first part is exact same !, just more pairs=(6x5+6x4+5x4) 
\subsection*{Answer2: indicator variable + linearity of expectation even though the $X_i$ not independent!!}
score 1 if slot 1,2 = good, 0 otherwise, and same for slot 2,3 etc ($\leftarrow$ the indicator variable). First slot 1,2 is good if (BM or MB) ie $\frac{8}{15}\frac{7}{14} + \frac{7}{15}\frac{
8}{14}$, same for slot 2,3 ... slot 14,15 so $14 ( \frac{8}{15}\frac{7}{14} + \frac{7}{15}\frac{
8}{14} ) = (2) \frac{56}{15} $ ($\leftarrow$ linearity of expectation even if not independent!!)

\section{ $2^n$ Tennis tournament with elimination}
\subsection*{Question: Prob second-best makes it to runner-up?}
$2^3=8$ players, better one always defeats next best. \textbf{Question}: prob second best is in final ?
\subsection*{Answer: Second best must be in lower half to meet the best in the final {\tiny (where he will be defeated of course)}}
\textbf{1 2 3 4} {\small 5 6 7 8}  ($\leftarrow$ second-best must occupy a lot in lower half to win it)= $\frac{\text{good slots for 2nd best}}{\text{possible slots for 2nd-best -- keep 1 for best}}=\frac{2^{n-1}}{2^{n}-1}$

\section{ $2^n$ tournament - prob of meeting at all}
\subsection*{Question: Chance of meeting in tournament}
\subsection*{Answer: enumerate}
\textbf{1 2} {\small 3 4 5 6 7 8} <= $\frac{1}{7}$ chance, and they meet with prob $1$,\textbf{1 3 or 1,4} {\small the rest} <= $\frac{1+1}{7}$ chance, and they meet with prob $\frac{1}{2},  \frac{1}{2}$ as they both have to win, finally bottom half: \textbf{1 5,6,7,8} with prob  $\frac{1+1+1+1}{7}$ and they meet with prob $\frac{1}{2x2x2x2}$ as they both have to win their matches and meet in final, so $\frac{1}{7} (1) + \frac{2}{7} \frac{1}{4} + \frac{4}{7} \frac{1}{16} = \frac{1}{4}$. Can prove general case by induction.

\section{Bernoulli trial - prob of given number}
\subsection*{Question: throw 100 fair coins, prob of 50 heads exactly?}
\subsection*{Answer: $\binom{100}{50} {(\frac{1}{2})}^{100} = \dfrac{100!}{50!50!}{(\frac{1}{2})}^{100}$ }
use sterling approximation to evaluate $n!=\sqrt{2\pi} {n}^{n+1/2}{e}^{-n}$

\section{1 6 with 1 dice, 2 6s with 2x6 dice, 3 6s with 3x6 dice?}
\subsection*{Answer: enumerate the probs for 1 6, 2 6s, }
prob of 1 6 exactly with 6 dice: (6 choose 1)(1 6 exactly) = $\binom{6}{1}\frac{1}{6} {\frac{5}{6}} ^{5}$ , prob of x 6s exactly with n dice: (n choose x)(x 6s exactly in n trials) = $\binom{n}{x} {\frac{1}{6}}^{x} {(\frac{5}{6})} ^{n-x}$

\section{3 way duel}
\subsection*{Question}
A,B,C fight duel. A shoots and wins with prob 0.3, B with prob 1, C with prob 0.5. Best strategy for A ?
\subsection*{Answer: enumerate the options}
if A shoots C , B will shoot A so no good - if A shoots B and misses then B will choose to shoot C since C better than A, then A gets to shoot B with his prob =0.3 - if A shoots B and wins, then C <-> A <-> C <-> A ... untill one dies, and the prob sequence is P(A eventual win)= P(A terminates the  sequence) = (C miss, A wins) + (C miss, A miss, C miss, A wins) + ... = (.5)(.3) + $(.5)^2 (.7)(.3)$ + $(.5)^3 (.7)^2 (.3) + ... = (.5)(.3) (\text{(geometric of .5 .7)} < 0.3$ so A should deliberately miss B and let B take care of C and then try hitting B with his prob 0.3

\section{Sample with and without replacement}
Urn A has 2R 1B, Urn B has 101 R,100 B. try and guess the right urn based on 2 draws from same urn: 1st draw, see the colour, choose to replace it or not replace it before drawign again form same urn. which is best ?
\subsection*{Answer: enumerate}
if 1st is red and urn 1, then 1B-1R and if urn2 100B -100R so 2nd draw from same urn is the same, so cannto tell which urn <-- put it back if red \\
if 1st is black, then if no replace, then the probs depends on the sequence of colours observed as follows:\\
\begin{tabular}{|c|c|c|c|}
	\hline 
	& prob it is urn A & prob it is urn B & biggest \\ 
	\hline 
	R*,R&(choose A)(R)(R) = $\frac{1}{2}\frac{2}{3}\frac{2}{3}$&$\frac{1}{2}\frac{101}{201}\frac{101}{201}$& urn A \\ 
	\hline 
	R*,B& $\frac{1}{2}\frac{2}{3}\frac{1}{3}$ & $\frac{1}{2}\frac{101}{201}\frac{100}{201}$ &  B \\ 
	\hline 
	B,R& $\frac{1}{2}\frac{1}{3}\frac{1}{1}$ & $\frac{100}{201}\frac{101}{200}\frac{1}{1}$ & A - note the no rep \\ 
	\hline 
	B,B& $\frac{1}{2}\frac{1}{3}0$ &  &  B - note the 0 in A\\ 
	\hline 
\end{tabular} 


\section{Ballot theorem - part 1 - vote counting}
\subsection*{Question suppose candidate A has more votes than B: a>b. Prob of having same number at least once at some point during the vote counting?}
\subsection*{Answer}
We know A does eventually win so ALL sequences with B starts first are good sequences, since there is eventual hit-and-overtake by A. All sequences with B first=$\frac{b}{a+b}$ . Now each sequence has a mirror : B <-> A, so there is same number of sequence starting with A with eventual equal tally (then A wins). \textbf{There are no other good sequences with A start else they would have a mirror and we have already counted all the mirrors.} so answer = $\frac{2b}{a+b}$

\section{Ballot theorem - part2 - Bernoulli sums}

\subsection*{Question: Bernoulli sums of 1,-1 - prob of no tie}
\subsection*{Answer: for even N=2n ..., prob a string x prob of a tie as per above}
Prob of x wins: $\binom{N=2n}{x}$ each string prob $\frac{1}{2^N}$ so $\binom{N=2n}{x}\frac{1}{2^N}$, and to use part 1 we need a win, so if x<=n -> loss -> use part 1 direct -> prob of a tie = $\frac{2x}{N}$, if x<=n -> loss -> convert to loss to use part 1 -> -> prob of a tie = $\frac{2(N-x)}{N}$ then add up the possibilities x =0..n $\binom{N=2n}{x}\frac{1}{2^N}\frac{2x}{N} $ and for x =n+1..N  $\binom{N=2n}{x}\frac{1}{2^N}\frac{2(N-x)}{N} $ then expand out and remember fact that $(1+1)^N=2^N=\sum{\binom{N}{x}}$


\section{Subway arrival}
\subsection*{Answer: train1 arrives 3:00 3:10.., train 2 arrives 3:02 3:12... - must arrive in good interval to catch1 before 2}

\section{prob(length Random Chord) > radius}
\subsection*{Answer: several possibilities}

\begin{tabular}{c|c}
	
\begin{tikzpicture} [scale=1]
\draw (0.4,0.4) circle (0.4 cm);
\draw (0.4,0.4) -- (0.4,0) -- (0,0.4) -- cycle;
\draw (0.4,0.4) -- (0.2,0.2);
\end{tikzpicture}
&
(1) chord distance / radius = $\frac{\sqrt{r^2-(r/2)^2}}{r}=\sqrt{3}/2$ \\
\end{tabular} 
\\(2) midpoint of chord less than altitude of regular hexagon = $\frac{\pi d^2}{\pi r^2} = {(\sqrt{3}/2)}^{2}=3/4$\\
(3) 2 random points on circle : given 1 point , second must fall either way of corresponding arc subtended by hexagon = 2 x 1/6 = 1/3 so bigger than radius = 1-1/3=2/3\\

\section{probability of meeting if each arrive between 6/7am and stay 5mns}
\subsection*{Answer: stretch of time between $y=x+1/12$ and $y=x-1/12$ with x=0..1}
because y can arrive either way of $ [x-1/12,x+1/12] $ with x 0..1 and constraint 0<y<1

\section{fixed number of urns with sampling - part1}
100 coins per urn, 1 fake coin per urn, draw 1 coin per urn, prob of detecting one false coin?
\subsection*{Answers}
probability of no false coins in all 100 urns with 1 false coins per urn $(1-\frac{1}{100})^{100}$, for n: $(1-\frac{1}{n})^n$, probability of no false coins in all n urns with m false coins $(1-\frac{m}{n})^{n} \rightarrow e^{-m}$

\section{multiple urns with sampling - part 2 -- POISSON DISTRIBUTION INTRODUCTION}
n coins per urn, m fake coins per urn, draw 1 coin from the n urns, prob of drawing exactly r false coins?
\subsection*{Answers}
r events each with prob $p=m/n$ so P= (n choose r) (prob of a string of r good and n-r bad) =\framebox{$\binom{n}{r}\frac{m}{n})^{r}(1-\frac{m}{n}))^{n-r}$} , then with n large and m and r fixed \framebox{$\binom{n}{r}\frac{m}{n})^{r}(1-\frac{m}{n}))^{n-r} = \frac{1}{r!} \frac{n(n-1)(n-r+1)}{n^r}m^r (1-\frac{m}{n})^{n} (1-\frac{m}{n})^{-r} = \frac{1}{r!} m^r e^{-m}$} which is the POISSON distribution \framebox{$P(r)=e^{-m}\frac{m^r}{r!}  $ r=0,1,2..} -- NOTE: large n -> means r can be very large,=$P(0)=e^{-m}$ 

\section{POISSON: Fixed number of colonies per plate with plates divided into small }
\subsection*{Question1: large number of plates, fixed average m=3 number of colonies per plate, how many plates with exactly m=3 colonies ? Question 2: let m=3 be large, what happens ?}

divide each plates into large number of small subplates, let p = prob = num colonies per sub-plate, then there on a plate $m = 3 = \text{(num sub plates)}\text{(prob per sub plate)} = np$ ie $np$ stays constant. Then if prob of $r$ colonies in 1 plate = $n$ sub-plates, with each colony = prob $p$ per sub-plate is Bernoulli experiment !

So $P(r)=\binom{n}{r} p^r (1-p)^{n-r}$ with $np=3=m \rightarrow p= \frac{m=3}{n}$ So $P(r)=\binom{n}{r} \frac{3}{n}^r (1-\frac{3}{n})^{n-r} \text{ with n } \rightarrow \infty$

See Question 26 ! POISSON distribution \framebox{$P(r)=e^{-m}\frac{m^r}{r!}  $ m=3 and r=0,1,2..}, Now set m=3 , and then r=3=exact number to answer the question.

Actually m=3 is the mean of the POISSON  distribution: $mean = \sum_{0}^{\infty} xP(x) = \sum_{r=0}^{r=\infty} rP(r) = \sum_{r=0}^{r=\infty}  r e^{-m}\frac{m^r}{r!} = \sum_{r=1}^{r=\infty} e^{-m}\frac{m^r}{r-1!}=m e^{-m} \sum_{r=1}^{r=\infty} \frac{{m}^{r-1}}{r-1!} = m e^{-m} \sum_{0}^{\infty} \frac{{m}^{r}}{r!} = m e^{-m}e^{m}=m$

Again : POISSON distribution \framebox{$P(r)=e^{-m}\frac{m^r}{r!}$} r=0,1,2.. . |Check out what happens when mean m $\rightarrow \infty$ :: Use sterling formula for r! with r=m $\rightarrow \infty$ ie $r!=\sqrt{2\pi}r^{r+\frac{1}{2}}e^{-m}$ so $P(r=m)=e^{-m}\frac{m^m}{m!}=e^{-m}\frac{m^m}{\sqrt{2\pi}m^{m+\frac{1}{2}}e^{-m}}=\frac{1}{\sqrt{2\pi}m}$
	
\section{POISSON: Sell m=20 cakes per round - Prob of selling even number of cakes}
Number of cakes sold per round = POISSON distribution = \framebox{$P(r)=e^{-m}\frac{m^r}{r!}$} and 'selling evens' means r=0,or 2 or 4 etc so of course answer = $\sum_{r=\text{even}}^{\infty} P(r)$ now use expansions of $e^{m}=\sum \frac{m^r}{r!}$ and $e^{-m}$ to eliminate odd terms and divide by 2, and with the $e^{-m}$ in front of each expansion you get $\frac{1}{2} (1+e^{-2m})$ and set m=20 in there for the answer.

\section{BIRTHDAY problem: smallest num. of people required in group for at least 2 of them in the group to have same birthdate with prob > 1/2}

find prob of no person having same birthdate in group, take complement.

If N=365 birthdates, first person in group has $N$ dates to choose from, second has $(N-1)$ dates to choose from ... r th person has   $(N-r+1)$ dates to choose from so $N(N-1)..(N-r+1)$ , to be divided by each person having N dates to choose from = $N^r$ so $\frac{N(N-1)..(N-r+1)}{N^r}$ with N=365 and try rs such that the prob < 1/2 so that 1-prob > 1/2

approximation: $e^{-x}=1-x$ x small and $\frac{N(N-1)..(N-r+1)}{N^r} = 1(1-\frac{1}{N}) (1-\frac{2}{N}) .. (1-\frac{r-1}{N})$ so $e^{O}e^{-\frac{1}{N}}e^{-\frac{2}{N}}..e^{-\frac{r-1}{N}} = e^{-\frac{{\sum_{0}^{r-1} i}}{N}}=e^{-\frac{r(r-1)}{2N}}$

\section{PERSONAL BIRTHDAY problem: how many people needed in group for one in group to have same birthday as you with prob > 1/2}

Note \textbf{difference}: need person in the group to have same birthday as you - NOT same birthday as each other (and different from you). this is the PERSONAL birthday problem.

say N=365 then Each person has $\frac{N-1}{N}$ chance NOT having your birthday, so in group of people : (not having yoru birthdays) (not having yoru birthday) .. (not having your birthday)=${(\frac{N-1}{N})}^{r}$ so 1 of them has your birthday = $1-{(\frac{N-1}{N})}^{r}$ , now find r by trial and error such that this > 1/2 for N=365

approximation : same thing, $e^{-x}=1-x$ for small x so $1-{(\frac{N-1}{N})}^{r}=1-({1-\frac{1}{N})}^{r}=1-e^{-\frac{r}{N}}$ then set = 1/2, use logs.

\section{BIRTHDAY problem: link the 2 above}

- each person in personal birthday problems gives 1 chance to match your personal birthday\\
- each person in group shared birthday problem has r-1 chances to have same birthdate as someone else -> $r(r-1)/2$ pairs\\
- with the approximation set $n = r(r-1)/2$ so that both probs are equal

\section{BIRTHDAY problem: whole group gets a holiday if one of them has a birthday. how large the group to have max number of working days = num people x num days worked ?}

- Note: group of 1 : only gets 1 holiday per year -> 1 person works 364 days -> 1x364 days\\
- Note: group of 2 : each gets 2 holiday per year -> 2 people work 2 x 363 days\\
- Note: group massive : every one gets holidays whole time-> massive x 0 days = 0 days worked\\
- so max is in between somewhere [2 .. massive] - but where ?\\
- Answer : everyone works same number days (days when there are no birthday), so E(num days worked) = (num people) x E(num days worked by 1 person) = (num people) (num days in year) prob(day is working day) = nN [(not a birthday for person1) (not a birthday for person2)...(not a birthday for person n)]= $nN { ( 1-\frac{1}{N})}^{n}$ and find which n gives max for N=365\\
- note more rigorous : expected (num days offered by first working day for n people) = n x  [(+1)xProb(+1)+(+0)xProb(+0)] = $n \left\{ (+1){ \left(1-\frac{1}{N} \right )}^n  + (+0) \left( 1-{(1-\frac{1}{N})}^n  \right) \right \}$\\
- so max $xp^x$, turns out x=N so $N^2{(1-1/N)}^N ==N^2 e^{-N/N}=N^2 e^{-1}$ so just $e^{-1}$ per worker per day with N workers and N days 

\section{RANDOM WALK ON INTEGER LINE - CLIFFHANGER: 1 step away from falling, prob $p$ of stepping right=away}
Note: possible paths: 1->0, 1->2->1->0, 1->2->1->2->1->0 etc\\
\textbf{FSA Answer}: FSA: prob (x=0|start=1) = prob (left)+prob (step right) prob (x=0|start=2)=$(1-p)P_1+pP_2$ but  $P_2$=prob (x=0|start=2)=P[all paths that go from 2 to 0 somehow]=P[all paths that go from 2 to 1 somehow AND all paths that go from 1 to 0 somehow] = P[all paths that go from 2 to 1 - somehow] x P [all paths that go from 1 to 0 - somehow] = $P_1 P_1$ (because markov and going form 1 to 0 somehow = going from 2 to 1 somehow) [Note: somehow means : in any way, maybe numerous steps like above] , anyway we get $P_1 = (1-p)+p {P_1}^2$ so $P_1$ is either $1$ or $\dfrac{1-p}{p}$ and P must be <1 hence =1 for $p<1/2$ and = $\dfrac{1-p}{p}$ for p>1/2\\
Note: for $p=2/3$ this is $1/2$ and $p=1/2$ this is $1$ - implication for casino : $P_1=1/2$ to avoid bankruptcy requires you to have 2/3 chance of winning... tons of paths can lead you to x=0 / bankruptcy \\
\textbf{FSA Answer general}: $P_2={P_1}^2$ so $P_n={P_1}^n$ so P(x=0|start m)= 1 or ${(\dfrac{1-p}{p})}^m$
\textbf{Martingale Answer}: use de Moivre martinagle (p/q) ....

\section{RANDOM WALK ON INTEGER LINE - GAMBLER RUIN: 2 players M,N play each other, winning odds: p(M win)=1/3<p(N win)=2/3, and initial fortune: M=1<N=2. they play +/-1 untill one bankrupt. what prob (M=ultimate winner) ? Answer: DOUBLE USAGE of the above}
- M starts with m=1 and N starts with n=2. Random walk on integer, and When x=0 , then M bankrupt, when x=m+n then M win\\
- from above with unlimited resource ie $n=\infty$ then x=0 ie bankrupt with prob $ {\frac{q}{p}}^m$\\
- but 2 ways to reach x= 0 = go bankrupt , (1) paths that always stay < m+n and (2) paths that reach m+n and then go back to 0. If let Q = paths that always stay < m+n, then if do reach m+n it is like a new start at x=m+n and we know already that these will go back to x=0 with prob $ {\frac{q}{p}}^{m+n} \leftarrow$ double usage of the above\\
- so basically $ \text{prob x= 0 with n=}\infty = {\frac{q}{p}}^{m} = Q + (1-Q) {\frac{q}{p}}^{m+n} \leftarrow \text{solve for Q}$ and P=prob of reaching m+n ever = 1 - prob stay below m+n = 1-Q \\
- answer = $P= \frac{1-{\frac{q}{p}}^{m}}{{1-\frac{q}{p}}^{m+n}} \text{ and with p=q=1/2 we get 0/0 and with l Hospital rule}=\frac{m}{m+n}$

\section{start with 20 \$ and want to reach \$40: 20\$ on evens in one go or bet 1\$ on even each time ? }
-Answer: bet 1\$ each time = random walk on integer line same as above m=start=20, want extra n=20, prob p = 18/38 and so q=1-p=20/38
- so use formula above for 1\$ at a time and compare with all in one go = p = 18/38

\section{PRINCIPLE - drop n points on a circle - then all symmetric ie length of (n+1) segments have identical distributions} 


\section{Mark sticks of length L with blue and red dots at the end , break sticks in 3 places, average length of blue-dot bit ? } 
Answer : equally distributed so L/3
\section{52 cards, how many cards till first ace ? } 
4 aces -> 5 segments of 0--48 -> 48/5 cards + ace
\section{German tanks 1..N, see n=60 so what is N? } 
Answer 1: 1 point dropped in 1..N -> your n creates 2 segments -> symmetry -> 2 segments of 1..59 + your obs -> 2x59+ the one you see = 2x59+1=119//
Answer 2: prob of sample containing the max = (1 x (N-1)) / $\binom{N}{2} = 2/N$ , max likelihood = make that prob biggest -> N=60 = as small as possible = what you see = n//
Answer 3: confidence interval = P(n you see>N/x)=some value ) eg P(n>2/3 N)=1/3 -> ???  

\section{German tanks 1..N, see 5 numbers n1...n5=60 so what is N? } 
Answer 1: 5 points -> 6 segments -> average length of (60-5)/5=11 -> N=60+11=71 

\section{break stick at random in 2 places, average length of smallest stick ? } 
Answer 1  break in 2 -> each segment average length L/2 by PRINCIPLE, now smallest = if break point left -> 1/2 x length of left = 1/2 x 1/2 = 1/4 \\
Answer 2: Smaller stick has length x which is uniformly distributed in the range [0,0.5]. The probability distribution of x is p(x) = 1/(b-a) = 1 / [0.5-0] = 2. Yes it is a pdf, $\int_{0}^{0.5} 2 dx = 1$. So E[x] = $\int x p(x) dx = \int_{0}^{0.5} x (2) dx = 2 \left[\frac{x^2}{2}\right]_{0}^{0.5} = 0.25 = 1/4$.
 
\section{break stick at random in 2 places, average length of ratio of lengths  ? } 
assume fell on right then x [1/2 .. 1] uniform dist 1/(b-a)=2 so $2\int_{1/2}^{1} \frac{1-x}{x}dx=2ln(2)-1$ 

\section{break stick at random in 2 places, average length of smallest middle, largest  ? } 
- from PRINCIPLE : break in 2 places -> 3 segments with uniform distribution -> average length of each segment = L/3\\
- however: so what ? we want SMALLEST and SMALLEST (3 uniforms) NOT uniform !! DIFFERENT from case above where the min was directly the distribution so was uniform \\
- Answer 0: complicated geometry methods with centroids \\
- Answer 1 FOR MIN: \\
$\bullet$ let z = min(x,y-x,1-y) = minimum length of the stick and P(z>a) = prob z > a \\
$\bullet$ P(z>a) = 2 * P(x>a, y-x>a, 1-y>a) = 2 P(x>a, a+x < y < 1-a) $\longleftarrow$ 2 because have to consider case x < y and y < x -- note MAX bit NOT entering yet thisis a general relation \\
$\bullet$ FIRST: a+x < y < 1-a $\rightarrow$ x < 1-2a so a < x < 1-2a $\rightarrow$  must have a < 1/3 ($Leftarrow$ obviously right , the min length will be less than 1/3 with 3 segments)\\
$\bullet$ SECOND: P(z>a) =  2 * 2 P(x>a, a+x < y < 1-a) $\rightarrow$ P(z>a) = 2 $\int_{a}^{1-2a} \int_{a+x}^{1-a} dy dx$ ($\leftarrow$ y first, then x) $\Rightarrow$ P(z>a) = 2 $\int_{a}^{1-2a} (1-2a-x) dx = 2 ( (1-2a) (1-2a-a) - 1/2 ((1-2a)^2 - a^2) = 2 (1-2a) (1-3a) - (1-3a) (1-a) = (1-3a) (2(1-2a)-(1-a) = (1-3a) (1-3a) = (1-3a)^2$ \\
$\bullet$ CONCLUSION: P(z>a) = $(1-3a)^2$ so P(z<a) = 1 - P(z>a) = $1-(1-3a)^2$ so p(a) = derivative of $1-(1-3a)^2 = 6 -18a = 6 (1-3a) \leftarrow$ btw : g'$\circ$f x f' = (-2)f x f' = (-2) (1-3a)(-3) = 6(1-3a) and E(a) = $ \int a p(a) da = \int_{0}^{1/3} 6(1-3a) da \leftarrow$ again remember 0 < a < 1/3 from above -- integration LIMITS !! \\
\\
- Now Answer 1 (ie calculus) FOR MAX: similar but stick to P(z<a)\\
$\bullet$ let z = max(x,y-x,1-y) \\
$\bullet$  Pr(z < a) = Pr(x < a, y-x < a, 1-y < a) = Pr( 1-2a < x < a, 1-a < y < x + a) $\leftarrow$ EXACT SAME AS ABOVE -- the MAX bit has not entered yet
$\bullet$ max bit now : ???

\section{your (P win) < P(lose) - How many $2n$ games to play to max your chances of winning? } 
$\bullet$ n+1 or more games in 2n to win so $ P(\text{win in 2n}) = P_{2n}=\sum \binom{2n}{k>= n+1} p^k q^{n-k}$\\
$\bullet$ what 2n that maximises that sum?\\
$\bullet$ if go on to n+2 and not stop at n, it must be because it is better for you, ie P(win in 2n+2) > P(win in 2n)\\
$\bullet$but 2n+2 games = 2n games + another 2\\
$\bullet$ in these 2 extra games your status changes in 2 cases only\\
$\bullet$ case 1: won n+1 in 2n but lose both 2 : then from winner ->  loser and you have reduced the prob by (lose x lose x n+1 wins in \\
$\bullet$ case 2: won only n in 2n but win both 2 : then now from loser -> winner\\
$\bullet$ and explicitly $P_{2n+2} = P_{2n} $ - P(2 losses) P(won n+1 | 2 losses) + P(2 wins)P(won n | 2 wins) and note P(won n+1 | 2 losses) = P(won n+1) becasue it is independent and same for the other and note ht - becasue that possibility REDUCES your prob of win and the other INCREASES it\\
$\bullet$ and so explicitly $P_{2n+2} = P_{2n} - q^2\binom{n+1}{2n}p^{n+1}q^{n-1} + p^2\binom{2n}{n}p^nq^n \longleftarrow$ note the MINUS to REDUCE the prob of winning and the PLUS \\
$\bullet$ and for optimal 2n we want P(2n) >= P(2n-2) and also P(2n) >= P(2n+2) \\
$\bullet$ and for this to happen the added bits need to be positive/negative\\
etc...

\section{MATCHING PROBLEM: 52 cards laid out, then another 52 cards laid out beneath - AVERAGE NUMBER of matches ?}
$\bullet$ COUNTING problem -> INDICATOR variables\\
$\bullet$ each card: $ 1 x P(\text{match}) + 0 x P(\text{not match}) = 1 x (1/n) + 0 x whatever = 1/n$ so n cards = n x 1/n = 1 (yes - don't care what comes before - E(sum of X) = sum of E(X) even if the Xs are dependent - so here what turns up IS dependent but the average is NOT dependent)

\section{MATCHING PROBLEM: secretary puts letters in envelopes at random : SAME PROBLEM - SAME AVERAGE}

\section{MATCHING PROBLEM: 52 cards laid out, then another 52 cards laid out beneath - PROBABILITY of num matches=r (not JUST the AVERAGE over all the probs as per above) ? this is a DERANGEMENT problem}
$\bullet$ P(r matches in n cards) = (prob 1st card matches) x (prob 2nd card matches) .. (prob rth card matches) x (choose r cards in n) x PROB(no match in n-r) $\rightarrow P(r|n)=\frac{1}{n} \frac{1}{n-1} \frac{1}{n-r+1} \binom{n}{r}P(0|n-r) = \frac{1}{r!} P(0|n-r) $\\
$\bullet$ P(0 matches in N cards) = num derangement of n cards / num card arrangements =  num derangement of n cards / n! \\
$\bullet$ DERANGEMENT = (1st card not in place) AND (2nd card not in place) AND .. (nth card not in place) = $ Num {T1}^c \cap {T2}^c \cap ... {Tn}^c = Num(\text{everything}) - Num ({T1} \cup {T2} \cup ... \cup {Tn}) =\lvert S\rvert-\sum_{i}\vert T_i\rvert+\sum_{i<j}\lvert T_i\cap T_j\rvert-\sum_{i<j<k}\lvert T_i\cap T_j\cap T_k\rvert+\cdots+(-1)^n\lvert T_1\cap\cdots\cap T_n\rvert = n!-\binom{n}{1}(n-1)!+\binom{n}{2}(n-2)!-\binom{n}{3}(n-3)!+\cdots+(-1)^n\binom{n}{n}(n-n)! = = n!-\frac{n!}{1!}+\frac{n!}{2!}-\frac{n!}{3!}+\cdots+(-1)^n\frac{n!}{n!}=n!\left[1-\frac{1}{1!}+\frac{1}{2!~}-\frac{1}{3!}+\cdots+(-1)^n\frac{1}{n!}\right].$\\
$\bullet$ so prob = $\frac{1}{r!} \text{Derangements} \frac{1}{n!} = \frac{1}{r!} \left[1-\frac{1}{1!}+\frac{1}{2!~}-\frac{1}{3!}+\cdots+(-1)^{n-r}\frac{1}{n-r!}\right]  \leftarrow $ goes to  1/e 1/r! 

\section{DECIDE IF MAXIMUM part 1: review N slips 1 by 1, decide if biggest or move on, if decision = actual biggest, then get it else get 0 - BEST STRATEGY ?}

$\bullet$ OPTIMAL POLICY : skip some , then select best so far : REASON: prob =k/n of having seen the max increases from 1/n to 1 but prob of winning decreases as we get near the end (and at the end it is just 1/n) so must be in the middle and any subsequent choice after your pick must be fails so whatever criterion you choose must be 'pick best so far and statement[k..n] is true if the true max is later' , so best policy is 'skip then pick' \\
$\bullet$ given optimal policy, the question is HOW MANY TO SKIP (s-1) before picking the sunsequent max so far (at k) : \\
$\bullet$ we want P(win with draw k = k is max so far) > P (win pick something after) \\
$\bullet$ P (win at k = k is the max) = 1/n \\
$\bullet$ P (max of the k-1 numbers seen so far is in the s-1 skipped -- else we have lost) = (s-1)(k-1) \\
$\bullet$ hence P (picking the maximum with the skipping strategy) = $ \\ \frac{1}{n}\sum_{k=s}^{n-1} \frac{s-1}{k-1}  $ \\
$\bullet$ hence to choose best we want P(win with draw s = s is max so far) > P ( win by picking something after s) \\
$\bullet$ ie $s/n > \frac{1}{n}\sum_{k=s}^{n-1} \frac{s-1}{k-1} = \frac{s}{n} \left[ \frac{1}{s} + \frac{1}{s+1} + ... + \frac{1}{n-1} \right] $ \\
$\bullet$ note limit: $ \frac{s}{n} \left[ \frac{1}{s} + \frac{1}{s+1} + ... + \frac{1}{n-1} \right] \equiv \frac{s}{n} ln (\frac{n}{s}) \rightarrow 1/e = s/n \text{ so cutoff } s \equiv n/e$

\section{DECIDE IF MAXIMUM part 2: same as before : review N slips 1 by 1, decide if biggest or move on, if decision = actual biggest, then get it else get 0 BUT the number on the slip is uniformly distributed !- BEST STRATEGY NOW ?}
$\bullet$ QUICK NOTE: before, did not know the dist of the dowry amounts. now we know the dist of the numbers. So say first number you get is 0.999, then prob (that this is > than next 99 numbers) = complement of (all number < 0.999) = 1 -\{ (2nd num < .999) and (3rd num < .999) .. \} = $1-(0.999)^{99}$ \\
$\bullet$ on this problem, work backward:
\begin{itemize}[leftmargin=0.5cm]
\item last draw : if biggest so far, then choose it (deterministic)
\item one draw before last : only one number can beat it (the last one), so if > 1/2 choose it, if not leave it, because if < 1/2 then prob (last one > 1/2) is > 1/2

\item 2 draws before last: then only next 2 draws can beat it. say current value is $x$ then  
\begin{itemize}[leftmargin=0.5cm]
\item both bigger : $(1-x)^2$
\item one bigger, one smaller : $x(1-x) + (1-x)x = 2x(1-x)$
\item both smaller: $x^2$
\item so prob of winning later is (CAREFUL) : (both bigger AND WE CHOOSE THE RIGHT ONE $\leftarrow$ ie prob  $\frac{1}{2}$ ) + (ONE BIGGER $\leftarrow$ so we are sure to choose it) = $\frac{1}{2}(1-x)^2 + 2x(1-x) $
\item so chance of winning later = $\frac{1}{2}(1-x)^2 + 2x(1-x) $
\item chance of winning now = now is bigger than later 2 = (bigger than next one) and (bigger than last one) = $x . x = x^2 \leftarrow$ uniform dist everywhere in this problem !
\item FINALLY : solve $x^2 = \frac{1}{2}(1-x)^2 + 2x(1-x) $ and choose candidate ie if bigger than solution
\end{itemize}

\item n draws before last: same reasoning , end up solving  $x^r = \binom{r}{1}x^{r-1}(1-x) + \frac{1}{2}\binom{r}{2}x^{r-2}(1-x)^2 + ... + \frac{1}{r}\binom{r}{r}x^{r-r}(1-x)^r$ 
\item approx: r big and x > 1/2 (always wait for this) then 1-x gets small so $x^r$ = first terma = $\binom{r}{1}x^{r-1}(1-x) \rightarrow x^r = rx^{r-1}(1-x) \rightarrow x = \frac{r}{r+1}$
\end{itemize}

\section{DOUBLE ACCURACY - measurement with noise}
$\bullet \quad D=A-B+d \quad S=A+B+s \quad A=\frac{1}{2}(D+S) - \frac{1}{2}(d+s)$ etc \\
$\bullet \quad$ average of $\frac{1}{2}(d+s)$ , variance of $\frac{1}{2}(d+s)$ = $\frac{1}{4} (\sigma_{s}^2 + \sigma_{d}^2) = \frac{1}{2} \sigma^2  $ 


\section{Probability that $X^2 +2 bX + c$ has real roots with b,c uniform over whole area}
$\bullet$ here the $b^2 -4ac$ is $4(b^2-c)$ so need $b^2-c  > 0$ so $b^2 > c$ \\
so choose b so that $b^2 > c$ \\
so choose b so that $b$ outside of $b^2 = c$ \\
$\bullet$ restrict b to [-B ... B ] and find area outside $b^2 = c$ \\
$\bullet$ then let B $\rightarrow \infty$
$\bullet$ area inside $b = \sqrt{c}$ integrate $ \int b = \int \sqrt{c} = \dfrac{ x^{\frac{3}{2}} } { \frac{3}{2}}$ with the +/- = $2 \frac{2}{3} x^{\frac{3}{2}} = \frac{4}{3} x^{\frac{3}{2}} = \frac{4}{3} B^{\frac{3}{2}} $\\
$\bullet$ whole area $4B^2$ so divide to get $\frac{1}{3 \sqrt{B}}$ \\
$\bullet$ NOTE : NOT SAME problem as $aX^2 +2 bX + c$ with a.b.c uniform because need divide by a and ratio of 2 uniforms not uniform !!

\section{2D random walk}
$\bullet$ in 2n moves, , just need n moves up ($\leftarrow$ and consequently n moves down) to return to origin, and this in both X and Y dimension
$\bullet$ so in 2n moves, P (X=0) = $\binom{2n}{n} {\frac{1}{2}}^n {\frac{1}{2}}^n = \binom{2n}{n} {\left( \frac{1}{2}\right) }^{2n}$ \\
$\bullet$ so in 2n x 2n  moves, P(return at 0) =  P (X=0)P (Y=0) = ${\left[  \binom{2n}{n} {\left( \frac{1}{2}\right) }^{2n} \right]  }^2$ \\
$\bullet$ sterling approximation $n!=\sqrt{2\pi} {n}^{n+1/2}{e}^{-n}$ so in 2n x 2n  moves, P(return at 0) = $ \dfrac{1}{\pi n}$ $\leftarrow$ straight calc \\
$\bullet$ so number of returns at origin = P(return after 2) (counts for 1 return) + P(return after 4) (counts for 1 return) + ... ($\leftarrow$ for ever) = $\frac{1}{\pi}\sum \frac{1}{n} = \infty$ \\
$\bullet$  so Q=0, P=1 since  Q = prob of termination in $\text{mean} = \mu = \sum x P^x Q$ and we know form 1D case $\mu = 1/Q$

\section{3D random walk}
$\bullet$ same principle P (X=0)P (Y=0) = ${\left[  \binom{2n}{n} {\left( \frac{1}{2}\right) }^{2n} \right]  }^3$ \\
$\bullet$ same thing in 2n x 2n x2n moves, P(return at 0) = $ \dfrac{1}{ {\left( \pi n \right) }^{3/2}}$ $\leftarrow$ straight calc $\leftarrow$ sum of this converges\\
$\bullet$ so this sum = mean number of returns = $ \sum \dfrac{1}{ {\left( \pi n \right) }^{3/2}} = \mu$ \\
$\bullet$ same as above with terminating sequence $\text{mean} = \mu = \sum x P^x Q$
$\bullet$ num of succsses + 1 failure = 1/Q -> Q = 1/(num success + 1) ($\leftarrow$ always < 1 !!) and $P=1-Q$


\section{Buffon SHORT needle of lengh $2l$ tossed over vertical lines $2a$ apart}
$\bullet$ compute distance of centre of needle to nearest parallel given angle with  parallel \\
$\bullet$ $l cos \theta$ , assume $\theta$ uniform over $ 0 .. \frac{\pi}{2}$ \\
$\bullet$ so average distance is $\dfrac{l \cos \theta}{a}$ with $\theta$ uniform over $ 0 .. \frac{\pi}{2}$ \\
$\bullet$ so $\int_{0}^{\dfrac{\pi}{2}} \left( \frac{1}{\frac{\pi}{2}}\dfrac{l \cos \theta}{a} d \theta\right)  = \frac{2l}{\pi a} \int_{0}^{\frac{\pi}{2}} \cos \theta d \theta = \frac{2l}{\pi a} [sin \theta]_0^{\frac{\pi}{2}} = \frac{2l}{\pi a} $
\end{spacing}

\section{Buffon SHORT needle of lengh $2l$ tossed over vertical and horizontal lines $2a$ apart}

\section{Buffon LONG needle of lengh $2l$ tossed over vertical lines $2a$ apart}

\end{multicols*}


\end{document}